\documentclass{article}
\usepackage{hyperref}
\usepackage{Style}

\nocite{*} % Comentar si quiero citar
%\addbibresource{bibliografia.bib} % Quitar el comentado si quiero usar bibliografia

\begin{document}

\begin{minipage}{2.5cm}
    \includegraphics[width=2cm]{imagen_puc.jpg}
\end{minipage}
\begin{minipage}{12cm}
    {\sc Pontificia Universidad Católica de Chile\\
    Facultad de Matemáticas\\
    Departamento de Matemática\\
    Profesor: Manuel Cabezas -- Ayudante: Benjamín Mateluna}
\end{minipage}
\vspace{2ex}

{\centerline{\bf Introducción a la Geometría - MAT1304}
\centerline{\bf Ayudantía 23}}
\centerline{\bf 11 de junio de 2026}

\vspace{1cm}
\noindent\textbf{Problema 1.} Demuestre que cualquier recta que pasa por $P=(-1,5)$ no puede ser 
tangente a la circunferencia de ecuación $x^{2}+y^{2}+4x-6y=-6$.

\vspace{5mm}
\noindent\textbf{Problema 2.} Pruebe que las cónicas $x^{2}+y^{2}+4x+6y-23=0$ y 
$x^{2}+y^{2}-8x-10y+25=0$ son tangentes.

\vspace{5mm}
\noindent\textbf{Problema 3.} Halle el lugar geométrico de los centros de las circunferencias que 
son tangentes a la recta $y=1$ y la circunferencia de ecuación $x^{2}+y^{2}=9$.

\vspace{5mm}
\noindent\textbf{Problema 4.} Encuentre el conjunto solución de la ecuación 
$5x^{2}+4xy+2y^{2}-24=0$.

\vspace{1cm}
\noindent\textbf{\underline{Ejercicios Propuestos:}}

\vspace{5mm}
\noindent\textbf{Extra 1.} Determine si la ecuación $4x^{2}+16x+4y^{2}-4y+17=0$ corresponde a
una cónica, un punto o el conjunto vacío. Si sucede lo primero, identifique que cónica es.

\vspace{5mm}
\noindent\textbf{Extra 2.} Describa el conjunto de soluciones de la ecuación 
$2x^{2}-5xy+3x-7y^{2}-6y+1=0$.

%\printbibliography % Quitar el comentado si quiero usar bibliografia

\end{document}